%%%%%%%%%%  scan idx 103  |  folio 92  %%%%%%%%%%

\underline{Corolario 3.14.}
Sea $R$ noetheriano, y sea $A$ un $R$-álgebra de tipo finito
que es lisa en todo punto salvo en un punto cerrado
$x\in\operatorname{Spec}(A)$. Entonces existen isomorfismos
\[
 D^1(A|R,A)\simeq D^1(A_x|R,A_x)
 \simeq D^1(\widehat A_x|R,\widehat A_x).
\]

\underline{Demostración.}
Como $A$ es lisa sobre $R$ fuera de $x$,
$D^1(A|R,A)$ es un $A$-módulo de tipo finito con soporte en
$\{x\}$. Por tanto,
\[
 D^1(A|R,A)=D^1(A|R,A)_x
 =\widehat{D^1(A|R,A)_x}.
\]
La afirmación se deduce ahora de 3.13.
\hfill$\square$

\underline{Observación 3.15.}
Sea $R$ un anillo topológico y $A$ un $R$-álgebra topológica.
Para todo $A$-módulo $E$, definamos
\[
 D^i_{\mathrm{top}}(A|R,E)
 =\varinjlim_\alpha
 D^i(A_\alpha|R_\alpha,A_\alpha\otimes_AE),
\]
donde $A_\alpha=A/J_\alpha$, $R_\alpha=R/I_\alpha$, y
$J_\alpha$ e $I_\alpha$ son abiertos en $A$ y $R$, respectivamente.
Así, $A$ es formalmente lisa como $R$-álgebra topológica precisamente
cuando $D^1_{\mathrm{top}}(A|R,E)=0$ para todo $A$-módulo $E$.

Sea ahora $X$ un esquema sobre un anillo $R$, y sea $\underline E$ un
$\mathcal O_X$-módulo cuasicoherente. Definimos haces
\[
 \underline D^i_{X|R}(\underline E)
 =\underline D^i(X|R,\underline E)
\]
sobre $X$ a partir de los prehaces sobre los abiertos afines
$U\subset X$ dados por
\[
 \bigl[\underline D^i_{X|R}(\underline E)\bigr](U)
 =D^i\bigl(\Gamma(U,\mathcal O_X)|R,
            \Gamma(U,\underline E)\bigr).
\]
Si $R$ es noetheriano y $X$ es de tipo finito sobre $R$, la
propiedad de localización 3.13 muestra que estos prehaces forman
canónicamente $\mathcal O_X$-módulos cuasicoherentes.

Cuando no haya posibilidad de confusión, abreviamos
\[
 \underline D^i_{X|R}
 =\underline D^i_{X|R}(\mathcal O_X).
\]
Todas las afirmaciones de 3.13 tienen formulaciones correspondientes
para los haces cuasicoherentes
$\underline D^i_{X|R}(\underline E)$. Dejamos esas formulaciones al
lector y registramos a continuación las propiedades más pertinentes
para nuestro propósito.

\underline{Corolario 3.15.}
Sea $R$ noetheriano y sea $X$ un $R$-esquema de tipo finito.

%%%%%%%%%%  scan idx 104  |  folio 93  %%%%%%%%%%

\begin{enumerate}[label=(\alph*)]
\item Para todo $\mathcal O_X$-módulo coherente $\underline E$,
$\underline D^i_{X|R}(\underline E)$ es también coherente para todo
$i$. Además,
\[
 \operatorname{Supp}\underline D^i_{X|R}(\underline E)
 \subset X^{\mathrm{Sing}}
 \qquad(i>0),
\]
donde
\[
 X^{\mathrm{Sing}}
 =\{x\in X\mid X\text{ no es liso sobre }R\text{ en }x\}.
\]

\item Si $X$ es plano sobre $R$, entonces para toda $R$-álgebra $S$
existe un isomorfismo canónico
\[
 \underline D^i_{X\otimes_RS\,|\,S}(\underline E)
 \xrightarrow{\ \sim\ }
 \underline D^i_{X|R}(\underline E)
\]
para todo $i$.

\item Supongamos que $X$ es una intersección completa relativa sobre
$R$. Para todo homomorfismo sobreyectivo de anillos
$S\longrightarrow R$, existe una sucesión exacta
\[
0\longrightarrow\underline D^1_{X|R}(\underline E)
\longrightarrow\underline D^1_{X|S}(\underline E)
\longrightarrow
\underline{\operatorname{Hom}}_{\mathcal O_X}
 \bigl((I/I^2)\otimes_R\mathcal O_X,\underline E\bigr)
\longrightarrow0,
\]
donde $I=\ker(S\longrightarrow R)$. Si, además, $R$ es local
con cuerpo residual $k$ y $X$ es plano sobre $R$, entonces
\[
0\longrightarrow
 \underline D^1_{X_0|k}(\mathcal O_{X_0})
\longrightarrow
 \underline D^1_{X|S}(\mathcal O_{X_0})
\longrightarrow
 D^1(R|S,k)\otimes_k\mathcal O_{X_0}
\longrightarrow0,
\]
donde $X_0=X\otimes_Rk$.
\end{enumerate}

\underline{Demostración.}
La parte (a) se deduce de la finitud y de las descripciones en grados
bajos de 3.12, mientras que (b) se deduce de 3.12(3).

Para (c), como $X$ es una intersección completa relativa sobre $R$,
tenemos $\underline D^i_{X|R}(\underline E)=0$ para todo $i\geq2$.
La sucesión de transitividad 3.12(2) proporciona, por tanto, una
sucesión exacta
\begin{multline*}
0\longrightarrow\underline D^0_{X|R}(\underline E)
\longrightarrow\underline D^0_{X|S}(\underline E)
\longrightarrow\underline D^0_{R|S}(\underline E)
\longrightarrow\underline D^1_{X|R}(\underline E)\\
\longrightarrow\underline D^1_{X|S}(\underline E)
\longrightarrow\underline D^1_{R|S}(\underline E)
\longrightarrow0,
\end{multline*}
donde los haces cuasicoherentes
$\underline D^i_{R|S}(\underline E)$ sobre $X$ se definen en todo
abierto afín $U\subset X$ mediante
\[
 \bigl[\underline D^i_{R|S}(\underline E)\bigr](U)
 =D^i(R|S,\underline E(U)).
\]
Como $S\longrightarrow R$ es sobreyectivo,
$\underline D^0_{R|S}(\underline E)=0$; y la descripción en grados
bajos de 3.12(8) da
\[
 \underline D^1_{R|S}(\underline E)
 =\underline{\operatorname{Hom}}_{\mathcal O_X}
  \bigl((I/I^2)\otimes_R\mathcal O_X,\underline E\bigr).
\]
Esto demuestra la primera sucesión exacta.

%%%%%%%%%%  scan idx 105  |  folio 94  %%%%%%%%%%

Si $X$ es plano sobre $R$, entonces, para todo
$\mathcal O_{X_0}$-módulo cuasicoherente $\underline E$, existen
isomorfismos canónicos
\[
 \underline D^i_{X_0|k}(\underline E)
 \xrightarrow{\ \sim\ }
 \underline D^i_{X|R}(\underline E)
\]
para todo $i$. En particular,
\[
 \underline D^1_{X_0|k}(\mathcal O_{X_0})
 \xrightarrow{\ \sim\ }
 \underline D^1_{X|R}(\mathcal O_{X_0}),
\]
de donde se deduce la segunda sucesión exacta.
\hfill$\square$

\underline{Observación 3.16.}
Los haces cuasicoherentes $\underline D^i_{X|R}(\underline E)$
definidos anteriormente deberían aparecer como los términos
$E_2^{0,i}$ de una sucesión espectral que relacione la teoría local
con la teoría global. L.~Illusie ha obtenido una teoría de este tipo
en \emph{Lecture Notes in Mathematics} 239. Su método utiliza
«álgebra homotópica», que generaliza los métodos de M.~André y
D.~Quillen, en lugar de «álgebra cohomológica». Parece que el
enfoque adoptado aquí puede globalizarse mediante la cohomología de
Grothendieck en la categoría de esquemas, dotada de una topología
apropiada que generalice la de la categoría afín.
